% Ch4 self-study (part 1) -- "I do / You do" ladder for Zumdahl 4.1-4.6.
% Format per ch03-selfstudy, extended: each YOUR TURN carries FOUR
% questions. Bare final answers visible in every edition (\selfcheck);
% full solutions key-only. Examples disjoint from ch04-notes/worksheets.
\documentclass{apchem-worksheet}

\apcunitword{Chapter}
\apcunit{4}
\apcunittitle{Types of Chemical Reactions and Solution Stoichiometry}
\apcdoctitle{Self-Study \textbullet\ \S4.1--4.6, I~do / You~do}
\apcsubtitle{Zumdahl \S4.1--4.6 \textbullet\ PDF pp.\ 169--190 \textbullet\ four YOUR TURN questions per ladder}

\begin{document}
\apctitleblock

\begin{apcnote}[How to use these notes --- read this first]
Each skill comes as a \textbf{ladder}: a fully worked example in the
solid-framed box, then \textbf{four} YOUR TURN questions in the dashed
box --- same skill, new numbers, no help.

\begin{enumerate}[leftmargin=*,itemsep=0.15em]
  \item \textbf{Work all four} before looking at anything; write real
        work, not fragments.
  \item Check against the gray \emph{check:} line. All four right on the
        first try earns the tick on the tracker (last page).
  \item Any misses: re-read the worked example, then redo the missed ones
        from scratch. Full solutions are in the teacher key.
\end{enumerate}
\end{apcnote}

% =====================================================================
\section*{\sffamily Ladder 1 \textbullet\ Water, and what dissolves in it}
\zsec{4.1--4.2}

Water is a \textbf{polar} molecule --- a bent frame with the oxygen end
partially negative and the hydrogen ends partially positive. Dissolving
an ionic solid is an electrical event: each ion is pried away and
surrounded by water molecules facing it charge-to-charge
(\term{hydration}).

What a dissolved substance does to conductivity sorts everything into
three bins:

\begin{center}
\renewcommand{\arraystretch}{1.2}
\begin{tabular}{@{}lll@{}}
\hline
\textbf{Bin} & \textbf{Who is in it} & \textbf{In solution} \\
\hline
strong electrolyte & soluble salts; strong acids and bases &
  ions only --- ionizes $\approx 100\%$ \\
weak electrolyte & weak acids and bases &
  mostly intact molecules, few ions \\
nonelectrolyte & molecular substances (sugars, alcohols) &
  molecules only; no ions \\
\hline
\end{tabular}
\end{center}

\begin{workedexample}[I do: sorting three solutes]
Classify \ce{NaNO3}, \ce{HC2H3O2} (acetic acid), and sucrose
(\ce{C12H22O11}) --- and say what is actually floating in each solution.

\textbf{\ce{NaNO3}} --- a soluble salt, so a \textbf{strong electrolyte}.
The solution contains \ce{Na+} and \ce{NO3-} ions and essentially no
intact \ce{NaNO3} units at all.

\textbf{\ce{HC2H3O2}} --- an acid, but not one of the six strong ones, so
a \textbf{weak electrolyte}. About 99\% of the molecules stay intact;
roughly 1\% ionize to \ce{H+} and \ce{C2H3O2-}. It conducts, but dimly.

\textbf{Sucrose} --- molecular, neither acid nor base, so a
\textbf{nonelectrolyte}. It dissolves beautifully (water hydrates its
polar \ce{-OH} groups) yet the solution carries no current --- dissolving
and ionizing are different events.
\end{workedexample}

\begin{yourturn}[YOUR TURN 1 --- four questions]
Classify each solute (strong / weak / non) and name the particles in
solution:
\begin{enumerate}[leftmargin=*,itemsep=0.5em]
  \item \ce{K2SO4} \workspace[1.2cm]
  \item \ce{HF} \workspace[1.2cm]
  \item methanol, \ce{CH3OH} \workspace[1.2cm]
  \item \ce{NH3} \workspace[1.2cm]
\end{enumerate}
\selfcheck{(a) strong --- \ce{K+} and \ce{SO4^2-} \quad (b) weak ---
mostly \ce{HF} molecules \quad (c) non --- molecules only \quad (d) weak
--- mostly \ce{NH3} molecules}
\keyonly{\begin{apcanswer}(a) Soluble salt $\to$ strong; two \ce{K+} and
one \ce{SO4^2-} per formula unit, no intact \ce{K2SO4}. (b) Acid but not
one of the six strong ones $\to$ weak; mostly intact \ce{HF}, a little
\ce{H+} and \ce{F-}. (c) An alcohol --- molecular, not an acid or base
$\to$ nonelectrolyte; hydrated \ce{CH3OH} molecules only. (d) Weak base
$\to$ weak electrolyte; mostly \ce{NH3} molecules with a little
\ce{NH4+} and \ce{OH-}.\end{apcanswer}}
\end{yourturn}

% =====================================================================
\section*{\sffamily Ladder 2 \textbullet\ Molarity, and the ions it
implies}
\zsec{4.3}

\[ M = \frac{\text{moles of solute}}{\text{litres of solution}} \]

For a strong electrolyte the label concentration is only the start ---
the \emph{ion} concentrations follow from the subscripts.

\begin{workedexample}[I do: a calcium nitrate solution, down to the ions]
\SI{8.20}{\gram} of \ce{Ca(NO3)2} ($M = \SI{164.10}{\gram\per\mole}$) is
dissolved in enough water to make \SI{250.0}{\milli\litre} of solution.
Find the molarity, then every ion concentration.

\textbf{Moles first:}
$\dfrac{8.20}{164.10} = \SI{0.0500}{\mole}$

\textbf{Molarity:}
$\dfrac{0.0500}{0.2500} = \SI{0.200}{\Molar}~\ce{Ca(NO3)2}$

\textbf{Ions, by subscript:} each formula unit releases one \ce{Ca^2+}
and \emph{two} \ce{NO3-}:
\[ [\ce{Ca^2+}] = \SI{0.200}{\Molar} \qquad
   [\ce{NO3-}] = 2 \times 0.200 = \SI{0.400}{\Molar} \]
There is no ``\ce{Ca(NO3)2} concentration'' in the beaker at all --- a
strong electrolyte exists in solution only as its ions.
\end{workedexample}

\begin{yourturn}[YOUR TURN 2 --- four questions]
\begin{enumerate}[leftmargin=*,itemsep=0.5em]
  \item Molarity when \SI{0.250}{\mole} of \ce{KBr} is dissolved to make
        \SI{750.0}{\milli\litre}: \workspace[1.4cm]
  \item Moles of \ce{NaOH} in \SI{125}{\milli\litre} of a
        \SI{0.400}{\Molar} solution: \workspace[1.4cm]
  \item Mass of glucose, \ce{C6H12O6}
        ($M = \SI{180.16}{\gram\per\mole}$), needed for
        \SI{500.0}{\milli\litre} of \SI{0.0250}{\Molar} solution:
        \workspace[1.8cm]
  \item $[\ce{Cl-}]$ in \SI{0.150}{\Molar} \ce{AlCl3}: \workspace[1.2cm]
\end{enumerate}
\selfcheck{(a) \SI{0.333}{\Molar} \quad (b) \SI{0.0500}{\mole} \quad
(c) \SI{2.25}{\gram} \quad (d) \SI{0.450}{\Molar}}
\keyonly{\begin{apcanswer}(a) $0.250/0.7500 = \mathbf{0.333}$~M.
(b) $0.125 \times 0.400 = \mathbf{0.0500}$~mol.
(c) $0.5000 \times 0.0250 = 0.0125$~mol $\times\, 180.16 =
\mathbf{2.25}$~g.
(d) Three \ce{Cl-} per \ce{AlCl3}: $3 \times 0.150 =
\mathbf{0.450}$~M.\end{apcanswer}}
\end{yourturn}

% =====================================================================
\section*{\sffamily Ladder 3 \textbullet\ Dilution}
\zsec{4.3}

Adding water changes the volume, not the amount of solute. ``Moles before
$=$ moles after'' is the entire theory:
\[ M_1 V_1 = M_2 V_2 \]

\begin{workedexample}[I do: bench acid from the stock bottle]
How would you prepare \SI{500.0}{\milli\litre} of \SI{1.00}{\Molar}
acetic acid from the \SI{17.4}{\Molar} glacial stock?

\textbf{Moles needed:} $0.5000 \times 1.00 = \SI{0.500}{\mole}$

\textbf{Stock volume carrying that many moles:}
\[ V_1 = \frac{M_2 V_2}{M_1} = \frac{(1.00)(500.0)}{17.4}
       = \SI{28.7}{\milli\litre} \]

\textbf{The procedure matters:} measure \SI{28.7}{\milli\litre} of stock,
then add water \emph{to a total volume} of \SI{500.0}{\milli\litre} in a
volumetric flask --- not ``add \SI{500.0}{\milli\litre} of water,''
which would overshoot the volume and undershoot the concentration.
\end{workedexample}

\begin{yourturn}[YOUR TURN 3 --- four questions]
\begin{enumerate}[leftmargin=*,itemsep=0.5em]
  \item Volume of \SI{6.00}{\Molar} \ce{NaOH} stock needed for
        \SI{250.0}{\milli\litre} of \SI{0.150}{\Molar} solution:
        \workspace[1.6cm]
  \item Final molarity when \SI{15.0}{\milli\litre} of
        \SI{4.00}{\Molar} \ce{KCl} is diluted to
        \SI{250.0}{\milli\litre}: \workspace[1.6cm]
  \item To what final volume must \SI{10.0}{\milli\litre} of
        \SI{2.50}{\Molar} \ce{HNO3} be diluted to give
        \SI{0.125}{\Molar}? \workspace[1.6cm]
  \item In every dilution, which quantity is unchanged, and why does that
        make $M_1V_1 = M_2V_2$ true? \workspace[1.6cm]
\end{enumerate}
\selfcheck{(a) \SI{6.25}{\milli\litre} \quad (b) \SI{0.240}{\Molar} \quad
(c) \SI{200.}{\milli\litre} \quad (d) moles of solute}
\keyonly{\begin{apcanswer}(a) $V_1 = (0.150)(250.0)/6.00 =
\mathbf{6.25}$~mL. (b) $M_2 = (4.00)(15.0)/250.0 = \mathbf{0.240}$~M.
(c) $V_2 = (2.50)(10.0)/0.125 = \mathbf{200.}$~mL.
(d) The \textbf{moles of solute} --- water adds volume but no solute, and
since $n = MV$ on both sides of the event, $M_1V_1 = M_2V_2$ is just
``$n$ before $=$ $n$ after.''\end{apcanswer}}
\end{yourturn}

% =====================================================================
\section*{\sffamily Ladder 4 \textbullet\ Predicting precipitates}
\zsec{4.4--4.5}

Mix two salt solutions and the possible new pairings are found by
\textbf{swapping partners}; whether either pairing actually forms a solid
is decided by the \textbf{solubility rules}:

\begin{enumerate}[leftmargin=*,itemsep=0.15em]
  \item Group 1A (\ce{Na+}, \ce{K+}, \ldots) and \ce{NH4+} salts:
        soluble. Nitrates: soluble. No exceptions worth memorizing.
  \item Chlorides, bromides, iodides: soluble \emph{except} with
        \ce{Ag+}, \ce{Pb^2+}, \ce{Hg2^2+}.
  \item Sulfates: soluble \emph{except} with \ce{Ba^2+}, \ce{Pb^2+},
        \ce{Ca^2+}, \ce{Sr^2+}.
  \item Hydroxides: insoluble \emph{except} group 1A (and \ce{Ba^2+},
        marginally).
  \item Carbonates, phosphates, chromates, sulfides: insoluble
        \emph{except} group 1A and \ce{NH4+}.
\end{enumerate}

\begin{workedexample}[I do: potassium carbonate meets copper(II) sulfate]
Will a precipitate form when \ce{K2CO3(aq)} and \ce{CuSO4(aq)} are mixed?

\textbf{Ions in the pot:} \ce{K+}, \ce{CO3^2-}, \ce{Cu^2+},
\ce{SO4^2-}.

\textbf{Swap partners --- the two candidates:} \ce{K2SO4} and
\ce{CuCO3}.

\textbf{Consult the rules:} \ce{K2SO4} --- group 1A salt, soluble; stays
dissolved. \ce{CuCO3} --- a carbonate, and copper is not group 1A or
ammonium: \textbf{insoluble}. A pale blue-green solid forms.

\[ \ce{K2CO3(aq) + CuSO4(aq) -> CuCO3(s) + K2SO4(aq)} \]

The phases \emph{are} the answer --- writing \ce{(s)} on the right
compound is the whole point of the exercise.
\end{workedexample}

\begin{yourturn}[YOUR TURN 4 --- four questions]
Precipitate or no reaction? Give the formula and phase of any solid.
\begin{enumerate}[leftmargin=*,itemsep=0.5em]
  \item \ce{NaBr(aq) + Pb(NO3)2(aq)} \workspace[1.4cm]
  \item \ce{NH4Cl(aq) + K2SO4(aq)} \workspace[1.4cm]
  \item \ce{Ni(NO3)2(aq) + NaOH(aq)} \workspace[1.4cm]
  \item \ce{Na3PO4(aq) + CaCl2(aq)} \workspace[1.4cm]
\end{enumerate}
\selfcheck{(a) \ce{PbBr2(s)} \quad (b) no reaction \quad
(c) \ce{Ni(OH)2(s)} \quad (d) \ce{Ca3(PO4)2(s)}}
\keyonly{\begin{apcanswer}(a) Bromide meets \ce{Pb^2+} --- one of the
three exceptions: \textbf{\ce{PbBr2(s)}}. (b) Candidates \ce{NH4+}/
\ce{SO4^2-} and \ce{K+}/\ce{Cl-} are all soluble combinations:
\textbf{no reaction} --- all four ions stay in solution. (c) Hydroxide
with a non-1A metal: \textbf{\ce{Ni(OH)2(s)}}. (d) Phosphate with
calcium: \textbf{\ce{Ca3(PO4)2(s)}} --- note the charge-balanced
formula, $3(2+)$ against $2(3-)$.\end{apcanswer}}
\end{yourturn}

% =====================================================================
\section*{\sffamily Ladder 5 \textbullet\ Molecular, complete ionic, net
ionic}
\zsec{4.6}

One reaction, three languages: the \textbf{molecular equation} keeps
formulas whole; the \textbf{complete ionic} shows every strong
electrolyte as its ions; the \textbf{net ionic} deletes the
\term{spectator ions} --- the ones appearing identically on both sides
--- and keeps only the chemistry.

\begin{workedexample}[I do: the carbonate precipitation, three ways]
For the mixture from Ladder 4:

\textbf{Molecular:}
\[ \ce{K2CO3(aq) + CuSO4(aq) -> CuCO3(s) + K2SO4(aq)} \]

\textbf{Complete ionic} --- every \ce{(aq)} strong electrolyte splits;
the solid does \emph{not}:
\[ \ce{2K+ + CO3^2- + Cu^2+ + SO4^2- -> CuCO3(s) + 2K+ + SO4^2-} \]

\textbf{Cross out what never changed} (\ce{2K+}, \ce{SO4^2-}):
\[ \boxed{\;\ce{Cu^2+(aq) + CO3^2-(aq) -> CuCO3(s)}\;} \]

The net ionic is the reusable fact: \emph{any} soluble copper(II) salt
plus \emph{any} soluble carbonate gives this same equation. Charge check:
$2+ + 2- = 0$ on both sides.
\end{workedexample}

\begin{yourturn}[YOUR TURN 5 --- four questions]
\begin{enumerate}[leftmargin=*,itemsep=0.5em]
  \item Net ionic equation for \ce{NaBr(aq) + Pb(NO3)2(aq)}:
        \workspace[1.6cm]
  \item Net ionic equation for \ce{Ni(NO3)2(aq) + NaOH(aq)}:
        \workspace[1.6cm]
  \item List the spectator ions in question (a): \workspace[1.2cm]
  \item Why does \ce{K+} never appear in any net ionic equation in this
        chapter? \workspace[1.6cm]
\end{enumerate}
\selfcheck{(a) \ce{Pb^2+ + 2Br- -> PbBr2(s)} \quad
(b) \ce{Ni^2+ + 2OH- -> Ni(OH)2(s)} \quad
(c) \ce{Na+} and \ce{NO3-} \quad
(d) every potassium salt is soluble, so \ce{K+} is always a spectator}
\keyonly{\begin{apcanswer}(a) \ce{Pb^2+(aq) + 2Br-(aq) -> PbBr2(s)} ---
the 2 is required for charge balance. (b) \ce{Ni^2+(aq) + 2OH-(aq) ->
Ni(OH)2(s)}. (c) \ce{Na+} and \ce{NO3-} appear unchanged on both sides.
(d) Group 1A salts are soluble without exception, so \ce{K+} always
enters and leaves as a free aqueous ion --- a permanent
spectator.\end{apcanswer}}
\end{yourturn}

% =====================================================================
\section*{\sffamily Mastery tracker}

Tick a row only if \textbf{all four} YOUR TURN questions were right on
the first attempt.

\renewcommand{\arraystretch}{1.35}
\noindent\begin{tabular}{@{}c p{6.6cm} c p{4.4cm}@{}}
\toprule
\textbf{First try?} & \textbf{Skill} & \textbf{Ladder} &
  \textbf{If not, re-read\ldots}\\
\midrule
$\square$ & electrolyte classification & 1 & the three-bin table \\
$\square$ & molarity and ion concentrations & 2 & subscripts multiply \\
$\square$ & dilution & 3 & moles before $=$ after \\
$\square$ & predicting precipitates & 4 & swap, then the rules \\
$\square$ & net ionic equations & 5 & split, cross out, check charge \\
\bottomrule
\end{tabular}

\begin{apcnote}[Scoring yourself honestly]
5/5: continue to the \S4.7--4.11 self-study --- it assumes exactly these
skills. 3--4: redo the missed ladders tomorrow, not today. 2 or fewer:
the usual root causes are (1) forgetting that subscripts multiply ion
concentrations, (2) applying a solubility rule without its exceptions,
and (3) splitting the \emph{solid} in a complete ionic equation ---
only dissolved strong electrolytes split.
\end{apcnote}

\end{document}
